\documentclass[10pt]{article}
\usepackage{amsmath}
\usepackage{amssymb}
\usepackage{hyperref}
\usepackage[svgnames]{xcolor} 
\usepackage{listings}
\usepackage[skins,listings,breakable]{tcolorbox}

\definecolor{cppbg}{HTML}{F8F9FA}     
\definecolor{boxtitlebg}{HTML}{2B2D31}
\definecolor{cppblue}{HTML}{00599C}   
\definecolor{cppgreen}{HTML}{138A07} 
\definecolor{cppred}{HTML}{B11515}  
\definecolor{cppnumber}{HTML}{888888} 

\lstdefinestyle{cppstyle}{
    language=C++,
    commentstyle=\color{cppgreen}\itshape,
    keywordstyle=\color{cppblue}\bfseries,
    numberstyle=\tiny\color{cppnumber},
stringstyle=\color{cppred},
basicstyle=\ttfamily\footnotesize,
breaklines=true,
    showstringspaces=false,
    tabsize=4,                    
    numbers=none                     
}

\newtcblisting{cppbox}[1]{
    enhanced,
    skin=bicolor,
    colback=cppbg,
    colbacktitle=boxtitlebg,
    coltitle=white,
    colframe=boxtitlebg,
    boxrule=0.5pt,
    arc=4pt,
    drop shadow=black!10!white,
    fonttitle=\bfseries\sffamily,
    title=#1,
    listing only,
    listing options={style=cppstyle, numbers=left, numbersep=10pt},
    left=15pt,
    top=4pt,
    bottom=4pt,
    breakable
}



\title{E.A.6.6 - Edge Coloring}
\author{}
\date{}

\begin{document}
\maketitle
\vspace{-6em}

\section{Modellazione}

Introduciamo le variabili
\[
x_{e,i}  \quad e \in E, i \in \{1, 2, 3\}
\]
$x_{e,i} = \top \iff $ l'arco $e$ è colorato con il colore $i$.

\paragraph{}

Sia $\mathcal{T}$ l'insieme di tutti i cicli di lunghezza 3 presenti nel grafo $G$. Un elemento di $\mathcal{T}$ è un set di tre archi $\{e_1, e_2, e_3\}$ che connettono tre vertici distinti (assumiamo no self-loops) $u, v, w \in V$:
\[
\mathcal{T} = \{ \{e_1, e_2, e_3\} \subseteq E \mid e_1=\{u,v\}, e_2=\{v,w\}, e_3=\{w,u\} \text{ con } u \neq v \neq w \}
\]

La formula $\Phi(G)$ è composta dalla congiunzione di:

\begin{enumerate}
    \item ogni arco deve avere almeno un colore
    \[
    \bigwedge_{e \in E} (x_{e,1} \lor x_{e,2} \lor x_{e,3})
    \]
    
    \item per ogni triangolo e per ogni colore, tre archi di un triangolo non possono essere dello stesso colore
    \[
    \bigwedge_{\{e_1, e_2, e_3\} \in \mathcal{T}} \bigwedge_{i=1}^{3} (\neg x_{e_1,i} \lor \neg x_{e_2,i} \lor \neg x_{e_3,i})
    \]
\end{enumerate}

Come per Schur's Lemma, non è necessario imporre che ogni arco abbia esattamente un colore. Se si assegnano più colori a un arco senza violare i vincoli sui triangoli, la soluzione è comunque valida (è sufficiente scegliere uno qualsiasi dei colori assegnati).

\section{Istanziazione}

\[
\begin{aligned}
E = \{ & \{G1, H\}, \{G1, E\}, \{H, E\}, \{H, I\}, \{H, A\}, \{I, A\}, \{I, B\}, \{A, E\}, \{A, S\}, \\
& \{E, D\}, \{S, D\}, \{S, B\}, \{S, C\}, \{B, J\}, \{B, C\}, \{B, G2\}, \{J, G2\}, \{D, C\}, \{C, G2\} \}
\end{aligned}
\]

\paragraph{}
L'insieme dei triangoli $\mathcal{T}$ presenti nel grafo è composto dai seguenti 7 cicli di lunghezza 3:
\[
\begin{aligned}
\mathcal{T} = \{ & \{\{G1, H\}, \{G1, E\}, \{H, E\}\}, \\
& \{\{H, I\}, \{H, A\}, \{I, A\}\}, \\
& \{\{H, E\}, \{H, A\}, \{A, E\}\}, \\
& \{\{B, J\}, \{J, G2\}, \{B, G2\}\}, \\
& \{\{B, C\}, \{C, G2\}, \{B, G2\}\}, \\
& \{\{S, D\}, \{D, C\}, \{S, C\}\}, \\
& \{\{S, B\}, \{B, C\}, \{S, C\}\} \}
\end{aligned}
\]

\paragraph{}
La formula proposizionale $\Phi(G_{1.1})$ in CNF è la congiunzione dei seguenti vincoli:

\paragraph{vincoli di colorazione per archi}
\begin{itemize}
    \item $(x_{\{G1, H\},1} \lor x_{\{G1, H\},2} \lor x_{\{G1, H\},3})$
    \item $(x_{\{G1, E\},1} \lor x_{\{G1, E\},2} \lor x_{\{G1, E\},3})$
    \item $(x_{\{H, E\},1} \lor x_{\{H, E\},2} \lor x_{\{H, E\},3})$
    \item $(x_{\{H, I\},1} \lor x_{\{H, I\},2} \lor x_{\{H, I\},3})$
    \item $(x_{\{H, A\},1} \lor x_{\{H, A\},2} \lor x_{\{H, A\},3})$
    \item $(x_{\{I, A\},1} \lor x_{\{I, A\},2} \lor x_{\{I, A\},3})$
    \item $(x_{\{I, B\},1} \lor x_{\{I, B\},2} \lor x_{\{I, B\},3})$
    \item $(x_{\{A, E\},1} \lor x_{\{A, E\},2} \lor x_{\{A, E\},3})$
    \item $(x_{\{A, S\},1} \lor x_{\{A, S\},2} \lor x_{\{A, S\},3})$
    \item $(x_{\{E, D\},1} \lor x_{\{E, D\},2} \lor x_{\{E, D\},3})$
    \item $(x_{\{S, D\},1} \lor x_{\{S, D\},2} \lor x_{\{S, D\},3})$
    \item $(x_{\{S, B\},1} \lor x_{\{S, B\},2} \lor x_{\{S, B\},3})$
    \item $(x_{\{S, C\},1} \lor x_{\{S, C\},2} \lor x_{\{S, C\},3})$
    \item $(x_{\{B, J\},1} \lor x_{\{B, J\},2} \lor x_{\{B, J\},3})$
    \item $(x_{\{B, C\},1} \lor x_{\{B, C\},2} \lor x_{\{B, C\},3})$
    \item $(x_{\{B, G2\},1} \lor x_{\{B, G2\},2} \lor x_{\{B, G2\},3})$
    \item $(x_{\{J, G2\},1} \lor x_{\{J, G2\},2} \lor x_{\{J, G2\},3})$
    \item $(x_{\{D, C\},1} \lor x_{\{D, C\},2} \lor x_{\{D, C\},3})$
    \item $(x_{\{C, G2\},1} \lor x_{\{C, G2\},2} \lor x_{\{C, G2\},3})$
\end{itemize}

\paragraph{vincoli sui triangoli}

\begin{itemize}
    \item \textbf{triangolo $\{\{G1, H\}, \{G1, E\}, \{H, E\}\}$:}
    \begin{itemize}
        \item $(\neg x_{\{G1, H\},1} \lor \neg x_{\{G1, E\},1} \lor \neg x_{\{H, E\},1})$
        \item $(\neg x_{\{G1, H\},2} \lor \neg x_{\{G1, E\},2} \lor \neg x_{\{H, E\},2})$
        \item $(\neg x_{\{G1, H\},3} \lor \neg x_{\{G1, E\},3} \lor \neg x_{\{H, E\},3})$
    \end{itemize}

    \item \textbf{triangolo $\{\{H, I\}, \{H, A\}, \{I, A\}\}$:}
    \begin{itemize}
        \item $(\neg x_{\{H, I\},1} \lor \neg x_{\{H, A\},1} \lor \neg x_{\{I, A\},1})$
        \item $(\neg x_{\{H, I\},2} \lor \neg x_{\{H, A\},2} \lor \neg x_{\{I, A\},2})$
        \item $(\neg x_{\{H, I\},3} \lor \neg x_{\{H, A\},3} \lor \neg x_{\{I, A\},3})$
    \end{itemize}

    \item \textbf{triangolo $\{\{H, E\}, \{H, A\}, \{A, E\}\}$:}
    \begin{itemize}
        \item $(\neg x_{\{H, E\},1} \lor \neg x_{\{H, A\},1} \lor \neg x_{\{A, E\},1})$
        \item $(\neg x_{\{H, E\},2} \lor \neg x_{\{H, A\},2} \lor \neg x_{\{A, E\},2})$
        \item $(\neg x_{\{H, E\},3} \lor \neg x_{\{H, A\},3} \lor \neg x_{\{A, E\},3})$
    \end{itemize}

    \item \textbf{triangolo $\{\{B, J\}, \{J, G2\}, \{B, G2\}\}$:}
    \begin{itemize}
        \item $(\neg x_{\{B, J\},1} \lor \neg x_{\{J, G2\},1} \lor \neg x_{\{B, G2\},1})$
        \item $(\neg x_{\{B, J\},2} \lor \neg x_{\{J, G2\},2} \lor \neg x_{\{B, G2\},2})$
        \item $(\neg x_{\{B, J\},3} \lor \neg x_{\{J, G2\},3} \lor \neg x_{\{B, G2\},3})$
    \end{itemize}

    \item \textbf{triangolo $\{\{B, C\}, \{C, G2\}, \{B, G2\}\}$:}
    \begin{itemize}
        \item $(\neg x_{\{B, C\},1} \lor \neg x_{\{C, G2\},1} \lor \neg x_{\{B, G2\},1})$
        \item $(\neg x_{\{B, C\},2} \lor \neg x_{\{C, G2\},2} \lor \neg x_{\{B, G2\},2})$
        \item $(\neg x_{\{B, C\},3} \lor \neg x_{\{C, G2\},3} \lor \neg x_{\{B, G2\},3})$
    \end{itemize}

    \item \textbf{triangolo $\{\{S, D\}, \{D, C\}, \{S, C\}\}$:}
    \begin{itemize}
        \item $(\neg x_{\{S, D\},1} \lor \neg x_{\{D, C\},1} \lor \neg x_{\{S, C\},1})$
        \item $(\neg x_{\{S, D\},2} \lor \neg x_{\{D, C\},2} \lor \neg x_{\{S, C\},2})$
        \item $(\neg x_{\{S, D\},3} \lor \neg x_{\{D, C\},3} \lor \neg x_{\{S, C\},3})$
    \end{itemize}

    \item \textbf{triangolo $\{\{S, B\}, \{B, C\}, \{S, C\}\}$:}
    \begin{itemize}
        \item $(\neg x_{\{S, B\},1} \lor \neg x_{\{B, C\},1} \lor \neg x_{\{S, C\},1})$
        \item $(\neg x_{\{S, B\},2} \lor \neg x_{\{B, C\},2} \lor \neg x_{\{S, C\},2})$
        \item $(\neg x_{\{S, B\},3} \lor \neg x_{\{B, C\},3} \lor \neg x_{\{S, C\},3})$
    \end{itemize}
\end{itemize}


\section{Codec}

\begin{cppbox}{}
#include "Encoder.hpp"
#include "Problem.hpp"
#include <string>
#include <vector>

class EdgeColoring : public Problem {
private:
    struct Triangle {
        int e1, e2, e3;
    };

    int numVertices;
    int numEdges;
    std::vector<std::pair<int, int>> edges;
    std::vector<Triangle> triangles;

    inline std::string vName(int e, int i) const {
        return "x_" + std::to_string(e) + "_" + std::to_string(i);
    }

public:
    // il costruttore riceve i nodi e la lista degli archi, e precalcola i triangoli
    EdgeColoring(int V, const std::vector<std::pair<int, int>>& E) 
        : numVertices(V), numEdges(E.size()), edges(E) {
        
        // matrice di adiacenza che mappa (u, v) all'indice dell'arco
        std::vector<std::vector<int>> adj(V, std::vector<int>(V, -1));
        for (int i = 0; i < numEdges; ++i) {
            int u = edges[i].first;
            int v = edges[i].second;
            adj[u][v] = i;
            adj[v][u] = i; // grafo non orientato
        }

        // troviamo tutti i triangolo
        for (int i = 0; i < V; ++i) {
            for (int j = i + 1; j < V; ++j) {
                if (adj[i][j] != -1) {
                    for (int k = j + 1; k < V; ++k) {
                        if (adj[i][k] != -1 && adj[j][k] != -1) {
                            triangles.push_back({adj[i][j], adj[j][k], adj[i][k]});
                        }
                    }
                }
            }
        }
    }

    void generateConstraints(Encoder& encoder, EncodingContext& ctx, 
                             int tId, int tCount) override {
        
        // ogni arco deve avere almeno un colore (1, 2 o 3)
        for (int e = tId; e < numEdges; e += tCount) {
            std::vector<int> edgeLits;
            for (int i = 1; i <= 3; ++i) {
                edgeLits.push_back(encoder.getVar(vName(e, i)));
            }
            encoder.writeDirectClause(edgeLits, ctx);
        }

        // per ogni triangolo e per ogni colore, 
        // tre archi non possono essere dello stesso colore
        // Partizionato per indice del triangolo
        for (int t = tId; t < triangles.size(); t += tCount) {
            int e1 = triangles[t].e1;
            int e2 = triangles[t].e2;
            int e3 = triangles[t].e3;

            for (int i = 1; i <= 3; ++i) {
                int lit1 = -encoder.getVar(vName(e1, i));
                int lit2 = -encoder.getVar(vName(e2, i));
                int lit3 = -encoder.getVar(vName(e3, i));
                
                encoder.writeDirectClause({lit1, lit2, lit3}, ctx);
            }
        }
    }
};
\end{cppbox}

\end{document}
